% !PRINTED XVI
\weberpage{XVI}
\inhaltpar{189}{Polygone}{667}
\inhaltpar{190}{Verzweigungspunkte und Verzweigungszahlen}{669}
\inhaltpar{191}{Polygonquotienten und Polygonklassen}{670}
\inhaltpar{192}{Polygonscharen}{672}
\inhaltpar{193}{Normalbasen}{676}
\inhaltpar{194}{Differentialquotienten}{679}
\inhaltpar{195}{Darstellung der Differentialquotienten durch Polygonquotienten}{682}
\inhaltpar{196}{Geschlecht des Körpers $\Omega$}{685}
\inhaltabschnitt{Siebenundzwanzigster Abschnitt.}{Algebraische und \gesperrt{Abel}sche Differentiale.}
\inhaltpar{197}{Differentiale in $\Omega$}{688}
\inhaltpar{198}{Die Polygonschar erster Gattung}{690}
\inhaltpar{199}{Der \gesperrt{Riemann-Roch}sche Satz}{695}
\inhaltpar{200}{Differentiale zweiter und dritter Gattung}{699}
\inhaltpar{201}{Die Residuen}{702}
\inhaltbuch{}{Tabellen.}
\inhaltzeile{I. Entwickelungen der sechzehn $\vartheta$-Quotienten (S.~88)}{711}
\inhaltzeile{II. Zweite Form der Entwickelung der sechzehn $\vartheta$-Quotienten (S.~91)}{713}
\inhaltzeile{III. Entwickelung der $\vartheta$-Quotienten in trigonometrischen Reihen (S.~92)}{716}
\inhaltzeile{IV. Entwickelungen der elliptischen Funktionen (S.~163)}{718}
\inhaltzeile{V. Entwickelung der Transzendenten zweiter Gattung (S.~164)}{720}
\inhaltzeile{VI. Verzeichnis von Klasseninvarianten (zum neunzehnten Abschnitt)}{721}
\inhaltzeile{\gesperrt{Alphabetisches Register}}{727}
